\chapter{Conclusion}
\label{ch:conclusion}

Au terme de ce travail, il convient de prendre du recul sur le chemin parcouru.
Plutôt que de récapituler chapitre par chapitre, je reviens d'abord sur les six
objectifs annoncés en introduction (section~\ref{sec:objectifs}) pour dire
honnêtement, pour chacun, ce qui a été réalisé, ce qui n'a été atteint que
partiellement, et ce qui reste ouvert. Je dégage ensuite les contributions que
je crois réellement transférables, puis les perspectives, avant une remarque de
transparence sur les outils employés.

\section{Bilan du travail réalisé}
\label{sec:concl_bilan}

L'ambition initiale était de construire une \emph{pipeline complète} perception
$\rightarrow$ planification $\rightarrow$ contrôle, modulaire et reproductible,
sur un bras low-cost SO-101 piloté via l'écosystème open-source LeRobot. Cette
pipeline existe, fonctionne de bout en bout, et chaque module est couvert par
des tests automatiques. C'est le résultat central. Mais le cahier des charges
fixait six objectifs plus précis, et l'honnêteté commande de les reprendre un à
un sans gommer les zones partielles.

\paragraph{Objectif 1 --- Architecture perception--planification--contrôle
multi-caméras (\emph{réalisé}).} L'architecture est en place et opérationnelle :
trois caméras (deux en stéréo eye-to-hand, une en eye-in-hand) alimentent un
module de perception qui produit une scène 3D dans le repère de base, un module
de planification de saisie qui en dérive des poses d'effecteur, et un module de
contrôle (cinématique inverse, génération de trajectoire, commande moteur) qui
les exécute. L'orchestration en boucle fermée est assurée par
\texttt{pipeline.py}. C'est l'objectif le plus complètement atteint, et celui
qui correspond le mieux à ma motivation réelle : disposer d'une chaîne complète,
lisible et modifiable, par opposition à une boîte noire end-to-end.

\paragraph{Objectif 2 --- Sélection de points de vue contre les occlusions
(\emph{partiel}).} Le système exploite bien deux régimes de perception
complémentaires : la stéréo eye-to-hand donne une vue globale de la scène, et la
caméra eye-in-hand \texttt{cam\_2}, amenée presque à la verticale de l'objet à
la pose d'approche, fournit un raffinement plus fiable que la stéréo oblique
(intersection rayon-plan, réorientation si l'objet est vu nettement allongé).
C'est une forme de \emph{vision active}, mais minimale : le robot ne planifie
pas activement \emph{où regarder} pour lever une occlusion donnée ; il bénéficie
d'un second point de vue par construction, sans raisonnement explicite sur
l'occlusion. Une sélection de point de vue à part entière reste une ouverture.

\paragraph{Objectif 3 --- Trajectoires évitant les obstacles
(\emph{partiel/limité}).} C'est l'objectif le plus en retrait par rapport à
l'ambition affichée, et il faut le dire clairement. Il n'y a pas de planificateur
de chemin complet (de type RRT) : l'évitement repose sur des mécanismes
géométriques plus simples mais effectifs --- zones d'exclusion autour de la base
du robot, bornes de l'espace de travail dans \texttt{scene.json}, et surtout
l'ancrage de la profondeur de prise sur le plan de la table plutôt que sur le $Z$
bruité de la stéréo. Ces garde-fous suffisent pour la tâche pick-and-place
mono-objet visée, mais ne constituent pas une planification de trajectoire sans
collision au sens fort. Le RRT reste une perspective.

\paragraph{Objectif 4 --- Replanification en boucle fermée
(\emph{réalisé}).} Le pipeline ne planifie pas une fois pour toutes : à la pose
d'approche, la caméra eye-in-hand corrige la position cible (correction
$XY$ par intersection rayon-plan, réorientation éventuelle), et l'IK est
recalculée. La phase de saisie comporte une boucle de tentatives avec
vérification après levée et un essai de rattrapage. La perception mise à jour en
cours de mouvement modifie donc bien le plan : la boucle fermée est réelle.

\paragraph{Objectif 5 --- Stratégies de saisie adaptées à la géométrie
(\emph{réalisé}).} La saisie top-down est implémentée et adaptée à l'objet :
l'ouverture de la pince est calculée à partir de la largeur mesurée
(\texttt{TopDownGrasp}), l'orientation du poignet aligne les mâchoires en travers
du petit axe de l'empreinte (yaw libre pour les objets debout à empreinte ronde,
aligné sur le grand axe pour les objets couchés), et la profondeur est ancrée sur
la mi-hauteur de l'objet posé sur la table~\cite{bohg_data-driven_2014}.
La prise latérale, en revanche, n'a pas été implémentée : une prise horizontale
autour d'un flanc vertical demande d'orienter l'approche hors du plan de tangage
du bras, ce que la sous-actionation à 5~DDL ne permet pas de façon fiable. Le
top-down s'impose donc comme l'approche robuste, au prix de borner l'ensemble des
objets effectivement saisissables.

\paragraph{Objectif 6 --- Évaluation expérimentale (\emph{partiel}).} C'est la
limite la plus importante à reconnaître. La chaîne complète est implémentée et
s'exécute de bout en bout, mais sa \emph{fiabilité n'a pas encore été mesurée} :
la campagne \emph{formelle} et chiffrée --- taux de réussite par classe d'objet,
nombre de replans, comparaison V1/V2/SmolVLA --- reste à mener avec le script
dédié, puis à journaliser.
\todoF{Insérer ici la synthèse chiffrée de la campagne formelle une fois
exécutée et journalisée (script \texttt{experiment\_campaign.py}). Ne reporter
aucune valeur tant que la campagne n'a pas tourné.}

En résumé, trois objectifs sont atteints (1, 4, 5), un l'est de façon partielle
mais opérationnelle (2), un l'est de façon limitée (3), et un dépend d'une
campagne encore à conduire (6). Ce bilan contrasté est, je crois, plus utile
qu'un satisfecit : il montre que la pipeline complète tient, et il situe
précisément ce qui resterait à approfondir.

\section{Contributions}
\label{sec:concl_contributions}

Au-delà du système lui-même, plusieurs éléments me paraissent réutilisables
au-delà de ce travail.

\begin{itemize}
  \item \textbf{Une pipeline modulaire reproductible et auto-testée.} Chaque
        étape (FK, IK, triangulation, planification de saisie) est validée par
        des tests de cohérence interne (par exemple le roundtrip FK$\rightarrow$IK
        à environ 0,8~mm, ou le self-test synthétique de triangulation sous le
        micromètre, $<10^{-4}$~mm), ce qui permet d'attribuer toute erreur résiduelle à la
        calibration ou à la détection 2D plutôt qu'au code géométrique.
  \item \textbf{Une calibration hand-eye robuste et une calibration stéréo
        conjointe.} Le solveur hand-eye traite l'ambiguïté planaire et rejette
        itérativement les poses aberrantes ; la calibration stéréo conjointe
        (déduction de $T_{\text{base}\rightarrow\text{cam}_1}$ depuis
        $T_{\text{base}\rightarrow\text{cam}_0}$ et $T_{\text{cam}_0\rightarrow
        \text{cam}_1}$) a fait passer la cohérence stéréo de 21{,}46 à 1{,}95~mm
        et rendu inutile la compensation manuelle du biais.
  \item \textbf{La distinction \texttt{color\_mode}
        chromatic/black/white/gray} dans le seuillage HSV, qui traite
        explicitement le fait que le noir et le blanc ne sont pas des
        couleurs au sens HSV~\cite{forsyth_computer_2012} --- extension non
        triviale du seuillage classique.
  \item \textbf{Le détecteur \texttt{HFDetector} fondé sur
        OWL-ViTv2}~\cite{minderer_owlv2_2023}, open-vocabulary et cohérent avec
        la même bibliothèque \texttt{transformers} que les politiques LeRobot,
        ce qui maintient une stack homogène.
  \item \textbf{La triangulation du sommet et l'orientation en repère base.}
        Estimer la hauteur en triangulant le point haut-centre des deux bboxes
        corrige la surestimation du proxy en pixels ; estimer l'orientation du
        grand axe directement dans le repère de base (projection rayon-plan +
        ACP 2D) rend l'angle indépendant de l'orientation des caméras.
  \item \textbf{L'ancrage de la profondeur de prise sur le plan de la table}
        plutôt que sur le $Z$ stéréo, motivé par l'analyse de propagation
        d'erreur ($\delta Z \propto Z^2$) : une décision géométrique simple qui
        évite la composante la plus bruitée de la mesure.
  \item \textbf{La saisie asservie au couple} : fermeture pas à pas qui lit
        \texttt{Present\_Load} et s'arrête au contact, avec vérification après
        levée --- une boucle de retour effective sur du matériel low-cost et une
        pince TPU compliante.
  \item \textbf{La caractérisation de la sous-actuation à 5~DDL}, qui explicite
        une limitation matérielle : le bras atteint de nombreuses orientations
        dans son plan de tangage mais ne peut pas pointer fiablement une approche
        hors de ce plan, ce qui borne rationnellement les stratégies de saisie
        (et écarte la prise latérale).
\end{itemize}

Je tiens à corriger ici une formulation du brouillon : la \emph{comparaison
empirique} entre approche modulaire et imitation learning (V1/V2/SmolVLA)
n'est pas une contribution réalisée mais une perspective. Le travail a mis en
place la stack permettant cette comparaison et l'a cadrée conceptuellement
(section~\ref{sec:concl_ia} ci-après et chapitre~\ref{ch:control_eval}) ; il ne
fournit pas encore de chiffres comparatifs.

\section{Perspectives}
\label{sec:concl_perspectives}

Ce travail ouvre plusieurs directions, qui prolongent directement les zones
restées partielles du bilan.

\begin{itemize}
  \item \textbf{Planification de trajectoire sans collision} (objectif 3) : un
        planificateur de chemin de type RRT~\cite{lavalle_planning_2006}
        remplacerait les zones d'exclusion géométriques par un évitement
        d'obstacles à part entière, condition nécessaire pour passer d'une scène
        mono-objet à un environnement réellement encombré.
  \item \textbf{Vision active avancée} (objectif 2) : raisonner explicitement sur
        les occlusions pour décider \emph{où regarder}, au lieu de bénéficier
        passivement d'un second point de vue.
  \item \textbf{Évaluation formelle et comparaison de paradigmes} (objectif 6) :
        conduire la campagne chiffrée puis comparer la pipeline modulaire à une
        politique d'imitation learning telle que
        SmolVLA~\cite{shukor_smolvla_2025} sur le même robot et les mêmes tâches.
  \item \textbf{Passage à un bras 6~ou 7~DDL} : lever la contrainte
        d'orientation analysée au chapitre~\ref{ch:discussion}
        (section~\ref{sec:disc_pistes}) et rendre fiable la prise latérale.
  \item \textbf{Couplage avec un modèle de langage}
        (\textit{LLM-as-planner}) pour traduire des instructions de haut niveau
        en séquences de saisies, et aborder des tâches multi-objets
        séquentielles (préparation d'un plateau, manipulation d'outils).
\end{itemize}

\section{Remarque finale sur l'utilisation des outils IA}
\label{sec:concl_ia}

Par souci de transparence scientifique, je précise l'usage que j'ai fait
d'outils d'intelligence artificielle au cours de ce travail. J'ai employé des
assistants de programmation et de rédaction (notamment Claude Code) pour
accélérer l'écriture du code et la mise en forme de ce mémoire. Ces outils ont
servi de support : génération de squelettes de fonctions, relecture, reformulation.

En revanche, les décisions techniques structurantes --- le choix d'une
architecture modulaire plutôt qu'end-to-end, la convention de repère, la
stratégie de calibration hand-eye et stéréo conjointe, l'ancrage de la profondeur
sur la table, la formulation de l'IK et le diagnostic de la sous-actuation, ou
encore l'identification empirique de la convention de fermeture à 90\textdegree{}
de la pince --- relèvent de ma réflexion personnelle. Chacune a été validée non
par confiance dans l'outil mais par les tests de cohérence interne décrits dans
ce mémoire et par l'expérimentation sur le robot réel : c'est le matériel et les
self-tests qui ont tranché, pas le générateur de texte.
\todoF{Préciser, si l'UNIGE le demande, une estimation honnête de la proportion
de code et de rédaction assistés par IA, et renvoyer aux bonnes pratiques
académiques en vigueur. Voix de Maxence.}
